\documentclass[11pt]{amsart}
\usepackage{amsmath,amssymb,amsthm,mathtools}
\usepackage[margin=1in]{geometry}
\usepackage{enumerate}
\usepackage{hyperref}

\newtheorem{theorem}{Theorem}[section]
\newtheorem{proposition}[theorem]{Proposition}
\newtheorem{lemma}[theorem]{Lemma}
\newtheorem{corollary}[theorem]{Corollary}
\theoremstyle{definition}
\newtheorem{definition}[theorem]{Definition}
\newtheorem{problem}[theorem]{Problem}
\newtheorem{missinglemma}[theorem]{Missing Lemma}
\theoremstyle{remark}
\newtheorem{remark}[theorem]{Remark}
\newtheorem{observation}[theorem]{Observation}

\newcommand{\RH}{\mathrm{RH}}
\newcommand{\bbR}{\mathbb{R}}
\newcommand{\bbC}{\mathbb{C}}
\newcommand{\bbN}{\mathbb{N}}
\newcommand{\bbZ}{\mathbb{Z}}
\newcommand{\bbQ}{\mathbb{Q}}
\newcommand{\Koff}{\mathcal{K}_{\mathrm{off}}}
\newcommand{\negind}{\operatorname{neg.ind}}
\newcommand{\Wtwo}{\mathcal{W}_2}
\newcommand{\dm}{dm_\infty}
\newcommand{\dnu}{d\nu}
\newcommand{\Tl}{T_\lambda}
\DeclareMathOperator{\Real}{Re}
\DeclareMathOperator{\Imag}{Im}
\DeclareMathOperator{\sig}{sig}
\newcommand{\La}{\Lambda}

\title[Two Closures and a Survivor]{Two closures and a survivor:\\
the amplification route to the Riemann Hypothesis\\
through almost-prime sieve bounds}

\author{David Alejandro Trejo Pizzo}
\thanks{Independent Researcher. \texttt{dtrejopizzo@gmail.com}}


\begin{document}
\nocite{bgIK,bgMV,bgFI,bgTenenbaum}

\begin{abstract}
In a companion paper \cite{P41} we argued that every attempted proof of
the Riemann Hypothesis (RH) within the Connes--Consani--Moscovici (CCM)
framework presents RH as the boundary point of a one-parameter family whose
interior is unconditionally controlled, that no one-dimensional propagation
principle reaches that boundary, and that exactly two logical forms survive
the resulting no-go: \emph{two-dimensional propagation} (Form~A) and
\emph{index amplification} (Form~B), with a cheap auxiliary line
(Form~C, an effective version of Bagchi's self-approximation criterion).
We now report the outcomes. \emph{Form~C closes:} the self-approximation
density does not degenerate as the comparison parameter $d\to0$; the scale
degenerates linearly, and the underlying amplifier is self-blocking---the
anchored density that would force RH and the zero-density ceiling that
bounds it under $\neg\RH$ carry the same exponent, so any improvement to one
improves the other. \emph{Form~A closes:} the only joint heat-time/arithmetic
trace $\mathcal{T}(t,x)$ whose two edges are unconditional theorems vanishes
identically in the interior, because the forward De~Bruijn--Newman flow
preserves reality of zeros; every coupling inequality is vacuous.
\emph{Form~B survives.} We prove that the negative index is strictly
superadditive under the Krein tensor product,
$\negind(Q_1\otimes Q_2)=p_1q_2+q_1p_2$, giving unconditional exponential
amplification $\kappa(\psi^{\otimes k})=(4m)^k/2$ on the off-critical
subspace $\Koff$; and we execute the decisive experiment of \cite{P41,P40}:
the two-variable Weil functional restricted to a thickened multiplicative
diagonal $h(x_1,x_2)=g(x_1x_2)\varphi(x_1/x_2)$ has (i) a spectral side
realising the Minkowski sum of zero pairs with the predicted
$(8,8)$ signature, and (ii) an \emph{autonomous} arithmetic side supported
on almost-primes, $\sum_N g(N)\sum_{d\mid N}\Lambda(d)\Lambda(N/d)\,
\varphi(d^2/N)$---a weighted transform of $\Lambda\star\Lambda$ computable
without knowledge of the zeros. The Hadamard barrier does not reappear
inside the functional. The entire remaining difficulty concentrates in one
statement, the \emph{sieve-to-index bridge}: that an unconditional sieve
upper bound on the almost-prime functional controls the negative-index
growth of the tensor powers. We state it precisely, explain why the parity
phenomenon does not obstruct it (one must \emph{bound}, not
\emph{distinguish}), and identify the wall it would meet if it required
asymptotics rather than bounds---the pair-correlation wall, which is
\emph{not} a positivity wall. No proof of RH is claimed.
\end{abstract}

\maketitle
\tableofcontents
\newpage

\section{Introduction}

\subsection{Where we stand}
The companion paper \cite{P41} reduced the search for a proof of RH, within
the CCM framework, to a structural trichotomy. After a programme that
mapped seven mutually reinforcing walls (MW-1 through MW-7; see
\cite{P41,P40}), we argued that the obstruction has a single logical form:
RH is always the boundary value of a one-parameter family whose interior is
an unconditional theorem, and no one-dimensional propagation principle
(limit, analyticity, parabolic backward uniqueness, monotonicity, ergodic
recurrence) carries the interior control to the boundary. By elimination,
two argument-forms remain: two-dimensional propagation (Form~A) and index
amplification (Form~B). A third, cheaper line---an effective version of the
Bagchi self-approximation criterion (Form~C)---was proposed as a low-risk
probe.

This paper reports what happened when each was pursued. The arithmetic of
the report is simple: Forms~A and~C close, with precise reasons that are
themselves informative; Form~B survives a first decisive experiment and
concentrates the entire remaining difficulty into one unconditional lemma.
The purpose of the paper is to record the two closures honestly and to
state the surviving problem sharply enough that it can be attacked or
refuted.

\subsection{What this paper proves, and does not}
We prove three closure/structure theorems
(Theorems~\ref{thm:formC-selfblock},~\ref{thm:formA-degeneration},
and the index calculations~\ref{thm:tensor-index}--\ref{thm:signature}),
all unconditional. We do \emph{not} prove RH, nor any conditional statement
that would follow from a weaker hypothesis than is already known. The
contribution is the localisation of the difficulty: after this paper, the
amplification route to RH rests on exactly one missing lemma
(Missing Lemma~\ref{ml:bridge}), of a type---a sieve upper bound on a
bilinear almost-prime functional---that lies in the one region where
unconditional methods are known to operate.

\subsection{Framework}\label{ssec:framework}
We use the notation of \cite{P39,P40,P41}. Briefly:
$\dm(s)=(2\pi)^{-2}|\Gamma(1/4+is/2)|^2\,ds$ is the archimedean spectral
measure; $\dnu=dm_{\mathrm{full}}-dm_{\mathrm{full,on}}\ge0$ is the
unconditionally non-negative discrepancy measure with
$\RH\iff\dnu=0$ \cite[~89]{P40}; the CCM trace criterion holds at a
single scale, $\RH\iff\Tl[\lambda_0]=0$ \cite[]{P40}; and the bridge
theorem \cite{P35} identifies the Weil quadratic form's negative index with
the negative index of the operator $H_C$ on the Pontryagin space $(K,Q)$,
$\kappa(Q)=\negind(H_C)=2m$, localised on the off-critical subspace $\Koff$
\cite[Thm~8.1]{P40}. Here $m$ is the number of off-critical zero
quadruples; $\RH\iff m=0$.

\section{Form C closes: the self-blocking amplifier}\label{sec:formC}

\subsection{The proposal and its mis-framing}
Bagchi's criterion makes RH equivalent to the strong recurrence of
$\zeta$ \cite[]{P40}. The self-approximation family
$\mathrm{SR}_d$---comparing $\zeta(s+i\tau)$ to $\zeta(s+id\tau)$ on a disk
in the strip with positive lower density of admissible $\tau$---is proved
unconditionally for every $d\ne0$ (Baker's theorem on linear forms in
logarithms for algebraic irrational $d$; extensions of Nakamura,
Pa\'nkowski, Garunk\v{s}tis covering all real $d\ne0$), while $d=0$ is
exactly RH. Form~C asked: as $d\to0$, how do the constants degrade?
\cite{P41} posed a dichotomy---exponential degradation (quantified MW-2,
closed) versus polynomial (a genuine tension worth exploiting).

The dichotomy was mis-framed, and seeing why is the content
\cite[]{P40}.

\begin{theorem}[The amplifier is self-blocking]\label{thm:formC-selfblock}
The lower density $\delta(d,\varepsilon)$ of admissible shifts does not
degenerate as $d\to0$: there is $d_0(\varepsilon)>0$ with
$\delta(d,\varepsilon)\ge\delta^*(\varepsilon)>0$ uniformly for
$0<|d|<d_0(\varepsilon)$. What degenerates is the scale: the first
guaranteed shift satisfies $T_0(d,\varepsilon)\asymp C(\varepsilon)/|d|$,
linearly. Moreover the Rouch\'e amplification that would convert one
off-critical zero into a positive-density family of off-critical zeros is
bounded, under $\neg\RH$, by an anchored density with the same exponent
$\kappa(\sigma)$ as the unconditional zero-density ceiling
$N(\sigma,T)=O(T^{\kappa(\sigma)})$ that contradicts it. The threshold
density needed to force RH and the ceiling that caps it move together under
any improvement of $N(\sigma,T)$.
\end{theorem}

The proof \cite[]{P40} isolates a diophantine-free regime
($|d|<d_0(\varepsilon)$, where the truncated frequency system carries no
linear relation and Baker is not invoked), shows the degradation is at worst
polynomial in $1/|d|$ where diophantine input does enter, and---decisively
---shows the Rouch\'e step consumes density not globally but in the
\emph{anchored window} $J_d=\{\tau:d\tau\approx\gamma_0\}$, where the slow
copy is frozen and the statement collapses to universality against a
\emph{forbidden} target (a shifted $\zeta$, which has a zero in the disk):
the admissible density there is provably zero by Hurwitz. The wall is
binary, not a rate.

\begin{corollary}
A mobile choice $d=d(T)\to0$ does not help: anchoring
($|d|T=O(1)$) freezes the slow copy and collapses $\mathrm{SR}_{d(T)}$ to
$\mathrm{SR}_0=\RH$; equidistribution ($|d|T\to\infty$) destroys the anchor
and yields at most $o(T)$ replicated zeros. The two regimes are
complementary and exhaustive \cite[Lem.~6.1]{P40}.
\end{corollary}

\begin{observation}
Form~C is a clean instance of MW-2 (the arithmetic propagation wall) made
quantitative. The missing lemma that would close it (localised universality
against a forbidden target in short intervals) is, on inspection, a
re-encoding of RH itself. We record it as such and do not pursue it.
\end{observation}

\section{Form A closes: the degeneration theorem}\label{sec:formA}

\subsection{The proposal}
Form~A sought a two-dimensional propagation principle. Let
$\mathcal{T}(t,x)$ be the CCM discrepancy trace of the finite-Euler
approximant indexed by $x$, evolved by the De~Bruijn--Newman (DBN) heat
flow to time $t$. The proposal \cite[Prob.~6.2]{P41} was that the two
edges---$\mathcal{T}=0$ for $t\ge\Lambda$ (the DBN constant) and reality of
the approximant's zeros for finite $x$---together with a coupling
inequality between $\partial_t\mathcal{T}$ and $\partial_x\mathcal{T}$
might propagate control to the corner $(t,x)=(0,\infty)=\RH$, even though
neither axis reaches it alone.

\subsection{The closure}
The coupling never gets a chance to exist.

\begin{theorem}[Degeneration]\label{thm:formA-degeneration}
The only definition of $\mathcal{T}(t,x)$ for which \emph{both} edges are
unconditional theorems is one that vanishes identically on the interior of
the quadrant. Consequently $\partial_t\mathcal{T}\equiv0$,
$\partial_x\mathcal{T}\equiv0$, and every parabolic, monotone, or
log-convex coupling inequality is vacuous.
\end{theorem}

\begin{proof}[Reason]
The forward DBN heat flow preserves reality of zeros: if $\Xi_t$ has only
real zeros at one time it has only real zeros for all later times
(de~Bruijn 1950 \cite{deBruijn50}). But ``reality of zeros'' is precisely
the property defining the finite-$x$ (arithmetic) edge. Hence the trace
that is zero on the arithmetic edge is forced to be zero throughout the
forward-invariant interior. Making both edges theorems and the interior
non-trivial are incompatible \cite[Thm~3.2]{P40}.
\end{proof}

The corner $(0,\infty)$ thus remains a pure discontinuity, whose removal is
the one-dimensional continuity wall already identified in \cite[\S5.1]{P41}.
Form~A reduces to that wall.

\subsection{Two byproducts}
The analysis \cite[]{P40} produced two unconditional byproducts of
independent interest. First, the only non-degenerate ledger---the Weil side
at $x=\infty$---admits a closed form for the first kernel block:
for $N(\lambda_0)=1$,
\[
B_{\lambda_0}(r)=\tfrac18\operatorname{sech}^{1/2}(r)\bigl(105\operatorname{sech}^4 r-78\operatorname{sech}^2 r+5\bigr),
\]
whose per-prime increments are positive for prime powers $p^m\le7$ and
negative for $p^m\ge8$, and whose prime series \emph{diverges} like
$(5\sqrt2/8)\log X$. Hence the trace identity of \cite[Thm~4.1]{P40}
is strictly formal: the CCM test function sits exactly on the boundary of
the Weil class (the poles of the weight at $\pm i/2$). Second, the analysis
corrected two prior errors in the corpus: a sign in the evolution equation
\cite[]{P40} and an omitted on-critical block \cite[]{P40},
the latter invalidating an earlier sign heuristic and reopening the sign of
$\partial_t T|_{0}$ under $\neg\RH$ as a declared gap.

\section{Form B survives: amplification of the negative index}\label{sec:formB}

\subsection{The Deligne skeleton, transplanted}
Deligne's proof over function fields \cite{Del74} is, in skeleton, not a
positivity argument: an integer invariant (the weight), an amplification
(a tensor power multiplies the defect), a uniform sublinear bound
(rationality: the poles do not move), and the conclusion that an integer
bounded by something sublinear in its own multiple must vanish. The CCM
programme already supplies the integer invariant: $\kappa(Q)=2m$, localised
on $\Koff$. Form~B asks for the amplification.

\subsection{The negative index is strictly superadditive}
The first half of the falsifier stated in \cite[\S7]{P41}---``perhaps
$\kappa$ is exactly additive under every available product''---is false.

\begin{theorem}[Tensor index]\label{thm:tensor-index}
For finite-dimensional Krein spaces with quadratic forms of signature
$(p_1,q_1)$ and $(p_2,q_2)$,
\[
\negind(Q_1\otimes Q_2)=p_1q_2+q_1p_2.
\]
This is strictly superadditive in the negative index whenever
$(p_1-1)q_2+(p_2-1)q_1>0$.
\end{theorem}

\begin{proof}
In a Sylvester basis diagonalising each $Q_i$ to $\pm1$, the product basis
diagonalises $Q_1\otimes Q_2$; a product vector is negative iff exactly one
factor is negative, giving $p_1q_2+q_1p_2$ negative directions.
\end{proof}

\begin{theorem}[Exponential amplification on $\Koff$]\label{thm:amplify}
With $m\ge1$ off-critical quadruples, $\psi=(\Koff,Q|_{\Koff})$ has
signature $(2m,2m)$, and under the Krein tensor product
\[
\kappa(\psi^{\otimes k})=\tfrac12(4m)^k\;\ge\;k\cdot2m=k\,\kappa(\psi)
\qquad(k\ge1),
\]
with equality only at $k=1$. Moreover $m=0\Rightarrow\kappa(\psi^{\otimes k})=0$
for all $k$: the invariant of the powers detects RH at every level.
\end{theorem}

This is ingredient~(1)--(2) of the skeleton, unconditional. The full tensor
product on $K$ (not its restriction to $\Koff$) is infinite-dimensional and
gives $\kappa=\infty$ already at $k=2$; the legitimate, finite invariant is
the \emph{relative} index on $\Koff$ \cite[\S3]{P40}.

\subsection{The sterility lemma}
Why is this not already a proof? Because amplification alone is sterile.

\begin{proposition}[Sterility]\label{prop:sterility}
In the abstract Krein-tensor setting, both the amplification (a) and any
bound on $\kappa(\psi^{\otimes k})$ (b) are computed from the same signature
data. No contradiction can be extracted from (a) and (b) together unless (b)
is supplied by a source independent of the signature.
\end{proposition}

In Deligne the independent source is cohomological rationality: the poles of
the $L$-function of the tensor family lie at prescribed places, with an
additive error uniform in $k$, guaranteed by the finite-dimensionality of
the $\ell$-adic cohomology. The question for Form~B is whether the CCM
framework possesses an analogous independent finiteness source. It does not
possess it for free---the tensor product must first be \emph{realised}
analytically, as a functional with its own explicit formula---and this is
the experiment we now execute.

\section{The decisive experiment: two-variable Weil with a diagonal kernel}\label{sec:experiment}

\subsection{Setup}
We must realise the tensor square $(\Koff,Q|_{\Koff})^{\otimes2}$ as a
spectral block of a functional whose arithmetic side is computable without
the positions of the zeros (Missing Lemma~5.2 of \cite[]{P40}).
The natural candidate is the two-variable Weil functional. Applied to a
product kernel $h(x_1,x_2)=f_1(x_1)f_2(x_2)$ it factors as
$\mathcal{W}\otimes\mathcal{W}$, with zero term
$\sum_{\rho_1,\rho_2}\hat f_1(\rho_1)\hat f_2(\rho_2)$. The eigenvalues of
$H_C\otimes I+I\otimes H_C$ are the sums $\rho_1+\rho_2$; to see them we
must \emph{couple} the variables.

\subsection{The thickened diagonal}
The naive diagonal $h(x_1,x_2)=g(x_1x_2)$ diverges: with $u=x_1x_2$,
$v=x_1$, the inner $v$-integral $\int_0^\infty v^{\rho_1-\rho_2-1}\,dv$ has
no decay. The well-defined object is the \emph{thickened} multiplicative
diagonal
\begin{equation}\label{eq:thick}
h(x_1,x_2)=g(x_1x_2)\,\varphi(x_1/x_2),\qquad
\varphi\in C_c^\infty(\bbR_+^*).
\end{equation}

\begin{theorem}[Spectral side; the Minkowski sum]\label{thm:signature}
With $u=x_1x_2$, $v=x_1/x_2$ (Jacobian $dx_1\,dx_2=du\,dv/2v$, exponent
$v^{(\rho_1-\rho_2)/2-1}$), the Mellin transform of \eqref{eq:thick} is
\[
\hat h(\rho_1,\rho_2)=\tfrac12\,\hat g\!\Bigl(\tfrac{\rho_1+\rho_2}{2}\Bigr)\,
\hat\varphi\!\Bigl(\tfrac{\rho_1-\rho_2}{2}\Bigr),
\]
so the zero term becomes $\sum_{\rho_1,\rho_2}\hat g(\tfrac{\rho_1+\rho_2}{2})
\hat\varphi(\tfrac{\rho_1-\rho_2}{2})$: the Minkowski sum of zero pairs,
weighted in their difference. For one quadruple
$\rho\in\{\tfrac12\pm\delta\pm i\gamma_0\}$, the sixteen pairs
$(\rho_i,\rho_j)$ organise under $\rho\mapsto1-\bar\rho$ into eight
fixed-point-free orbits---eight hyperbolic planes---so the
$16$-dimensional block has signature exactly $(8,8)$, matching
$\negind(Q|_{\Koff}^{\otimes2})=8$ of Theorem~\ref{thm:tensor-index}, with
no degeneration.
\end{theorem}

\begin{remark}[$\varphi$ carries the index]
A finding not anticipated in \cite{P41,P40}: the thickening weight
$\hat\varphi$ is not a mere regulariser. Of the eight negative directions,
only two sit at the real sums $\rho_i+\rho_j-1=\pm2\delta$; the other six
live at purely imaginary sums, where the off-criticality $\Real=\tfrac12\pm\delta$
is registered \emph{only} by $\hat\varphi$ through $\Real\bigl(\tfrac{\rho_i-\rho_j}{2}\bigr)=\pm\delta$.
The difference variable is a genuine carrier of the negative index, not an
artefact of regularisation.
\end{remark}

\subsection{The arithmetic side is autonomous}
This is the decisive point. The Hadamard barrier MW-7 would reappear if the
mixed zeros$\times$primes terms of the doubled functional required the
positions of the zeros. They do not.

\begin{theorem}[Mixed terms reabsorb]\label{thm:mixed}
Applied slice by slice---fixing the first variable at a prime power $n$ and
using the one-variable explicit formula on $x_2\mapsto h(n,x_2)\in
C_c^\infty$---the apparent zeros$\times$primes terms of the two-variable
functional collapse to arithmetic$\times$arithmetic terms. MW-7 does not
appear inside $\Wtwo$ \cite[Thm~4.4]{P40}.
\end{theorem}

\begin{theorem}[Almost-prime functional]\label{thm:almostprime}
For the thickened diagonal \eqref{eq:thick}, the primes$\times$primes block
of $\Wtwo$ is
\begin{equation}\label{eq:almostprime}
\sum_{N\ge1} g(N)\sum_{d\mid N}\La(d)\,\La(N/d)\,\varphi\!\bigl(d^2/N\bigr),
\end{equation}
a $\varphi$-weighted transform of the Dirichlet convolution
$\La\star\La$---the almost-prime part of Selberg's
$\La_2(N)=\La(N)\log N+(\La\star\La)(N)$---together with three
finitely-supported terms with $g,\varphi$ interchanged. It is computable
from the primes alone, without the positions of the zeros
\cite[Thm~4.2]{P40}.
\end{theorem}

Theorems~\ref{thm:signature}--\ref{thm:almostprime} together verify the two
conditions of Missing Lemma~5.2 of \cite[]{P40}: the doubled Weil
functional $\Wtwo$ realises the tensor-square spectral block, and its
arithmetic side is autonomous, living on the almost-prime semigroup. The
candidate for the analytic realisation $X_2$ is thus
$\Wtwo$ restricted to thickened diagonals, with
$Q_2(h,h')=\Wtwo(h\star\tilde h')$.

\section{What we want to prove: the sieve-to-index bridge}\label{sec:bridge}

\subsection{The honest accounting}
The amplification (Theorem~\ref{thm:amplify}) and the realisation
(\S\ref{sec:experiment}) are in hand. By the sterility lemma
(Proposition~\ref{prop:sterility}), the contradiction requires an
\emph{independent} bound on the negative-index growth of the tensor powers.
The arithmetic side \eqref{eq:almostprime} is exactly such a potential
source: it is a functional on almost-primes, supplied by the primes, not by
the zeros. The remaining problem is to convert a bound on this functional
into a bound on the index.

\begin{missinglemma}[The sieve-to-index bridge]\label{ml:bridge}
Let $X_k$ be the $k$-fold thickened-diagonal Weil functional realising
$(\Koff,Q|_{\Koff})^{\otimes k}$, with arithmetic side the order-$k$
almost-prime functional generalising \eqref{eq:almostprime}. Prove an
unconditional upper bound of \emph{sieve type}
\[
\kappa\bigl(Q_k^{\mathrm{rel}}\bigr)\;\le\;f(k),\qquad f(k)/k\to0,
\]
on the relative negative index of $Q_k$ on the off-critical block, with
$f(k)$ controlled by the almost-prime functional---the analogue of
``the poles do not move,'' here a uniform sieve bound rather than a
positivity.
\end{missinglemma}

If proved, Missing Lemma~\ref{ml:bridge} closes the argument: combined with
the unconditional lower bound $\kappa(\psi^{\otimes k})\ge2mk$
(Theorem~\ref{thm:amplify}, in its robust sublinear form), it forces
$2mk\le f(k)$ for all $k$, hence $m=0$, i.e.\ RH. We claim no proof of it.

\subsection{Why parity does not obstruct it}
The almost-prime location of \eqref{eq:almostprime} is not an accident, and
it is the reason for cautious optimism. The parity phenomenon---the
obstruction that closed the $\omega$-class programme (Arc~A of \cite{P40})
and that limits all sieve methods---forbids \emph{distinguishing}
$\omega(n)\bmod2$, equivalently forbids sieves from producing
\emph{asymptotics} for sequences detected by sign. But Missing
Lemma~\ref{ml:bridge} asks for an \emph{upper bound} on a bilinear
almost-prime functional, not an asymptotic and not a sign discrimination.
Upper bounds on $\La\star\La$-type sums over almost-primes are exactly what
the Selberg sieve, Vaughan's identity, and the Heath-Brown decomposition
deliver unconditionally. The functional \eqref{eq:almostprime} lives in the
one region where the methods operate.

\begin{remark}[The wall it would meet, if any]
We state the risk plainly. If the bridge required not an upper bound but an
\emph{asymptotic} evaluation of \eqref{eq:almostprime}---if the index growth
were governed by the main term of the almost-prime functional rather than by
a majorant---then the relevant obstruction is the pair-correlation/twin-type
barrier (Montgomery--Goldston), the difficulty of evaluating
$\sum\La(n)\La(n+h)$-type correlations. This is a real wall, but it is
\emph{not} the positivity wall MW-1/MW-4: it is an evaluation barrier, not a
sign barrier. No route in the programme has previously terminated against
it. Whether Missing Lemma~\ref{ml:bridge} needs a bound (passable) or an
asymptotic (the pair-correlation wall) is the first question to settle, and
it is a finite analytic question about \eqref{eq:almostprime}.
\end{remark}

\subsection{The shape of a proof, and its falsifiers}
A proof of Missing Lemma~\ref{ml:bridge} would proceed in three steps, each
falsifiable:
\begin{enumerate}[\rm(1)]
\item \emph{Index from functional.} Express $\kappa(Q_k^{\mathrm{rel}})$ as
the count of negative eigenvalues of the bilinear form
$Q_k(h,h')=X_k(h\star\tilde h')$ on the thickened-diagonal test space, and
bound it by the number of sign changes of the arithmetic kernel
\eqref{eq:almostprime}. \emph{Falsifier:} if the negative index is not
controlled by the kernel's sign changes (e.g.\ if the archimedean block
dominates), the arithmetic side is irrelevant and the route is dead.
\item \emph{Sieve majorant.} Bound the almost-prime functional by a Selberg
majorant uniform in $k$. \emph{Falsifier:} if the only available bound grows
like $k$ or faster (not $o(k)$), the sterility lemma is not overcome.
\item \emph{Uniformity.} Verify $f(k)/k\to0$ with the implied constant
independent of $m$. \emph{Falsifier:} if uniformity in $m$ fails, the bound
cannot kill all $m\ge1$ simultaneously.
\end{enumerate}
Step~(1) is the immediate next target; it is finite and self-contained.

\section{Assessment}\label{sec:assessment}

We have closed two of the three surviving forms, with reasons that sharpen
the map: Form~C is MW-2 quantified into a self-blocking amplifier
(Theorem~\ref{thm:formC-selfblock}); Form~A is a degeneration forced by the
forward heat flow (Theorem~\ref{thm:formA-degeneration}). The survivor,
Form~B, has cleared its first decisive test: the negative index amplifies
exponentially and unconditionally (Theorem~\ref{thm:amplify}), and the
tensor square is realised by a two-variable Weil functional whose
arithmetic side is autonomous and supported on almost-primes
(Theorems~\ref{thm:signature}--\ref{thm:almostprime}), with the Hadamard
barrier absent from the functional.

The entire weight now rests on one unconditional statement, Missing
Lemma~\ref{ml:bridge}, of a kind that lives where sieve methods work, and
whose only foreseeable wall is an evaluation barrier distinct from every
positivity wall the programme has met. We make no prediction about whether
it is provable. We claim only that the search for the missing idea has
narrowed to a single, finite, falsifiable analytic question: does an
unconditional sieve bound on the almost-prime functional
\eqref{eq:almostprime} control the negative-index growth of its tensor
powers? That narrowing---not any claim of progress toward RH itself---is
the contribution of this paper.

\begin{thebibliography}{99}

\bibitem{CCM24} A.~Connes, C.~Consani, H.~Moscovici,
\emph{Spectral realisation of zeta zeros and the Riemann--Roch programme}
(2024).

\bibitem{Del74} P.~Deligne, \emph{La conjecture de Weil. I},
Publ. Math. IH\'ES \textbf{43} (1974), 273--307.

\bibitem{deBruijn50} N.~G.~de~Bruijn, \emph{The roots of trigonometric
integrals}, Duke Math. J. \textbf{17} (1950), 197--226.

\bibitem{Weil52} A.~Weil, \emph{Sur les ``formules explicites'' de la
th\'eorie des nombres premiers}, Comm. S\'em. Math. Univ. Lund (1952),
252--265.

\bibitem{P35} D.~A.~Trejo Pizzo, \emph{Pontryagin spectral index of the Weil form and the Connes scaling operator: a bridge theorem for the Riemann Hypothesis}, preprint, 2026.

\bibitem{P39} D.~A.~Trejo Pizzo, \emph{A trace and measure criterion for the Riemann Hypothesis from the zeta spectral triple}, preprint, 2026.

\bibitem{P40} D.~A.~Trejo Pizzo, \emph{The localized Weil form and its Kre\u\i n--Pontryagin index: localization, obstruction, and rigidity}, preprint, 2026.

\bibitem{P41} D.~A.~Trejo Pizzo, \emph{The exhaustion of one-dimensional
propagation: a structural no-go for the Riemann Hypothesis, and two
directions that survive it} (programme paper P41, 2026).

\bibitem{bgIK} H.~Iwaniec and E.~Kowalski, \emph{Analytic Number Theory}, Amer.\ Math.\ Soc.\ Colloq.\ Publ.\ \textbf{53}, AMS, 2004.
\bibitem{bgMV} H.~L.~Montgomery and R.~C.~Vaughan, \emph{Multiplicative Number Theory I: Classical Theory}, Cambridge Univ.\ Press, 2007.
\bibitem{bgFI} J.~Friedlander and H.~Iwaniec, \emph{Opera de Cribro}, Amer.\ Math.\ Soc.\ Colloq.\ Publ.\ \textbf{57}, AMS, 2010.
\bibitem{bgTenenbaum} G.~Tenenbaum, \emph{Introduction to Analytic and Probabilistic Number Theory}, 3rd ed., Grad.\ Stud.\ Math.\ \textbf{163}, AMS, 2015.
\end{thebibliography}

\end{document}
